\section{Budget Matching Algorithm}
To efficiently determine the threshold $\tau_{\bar{N}_\mathcal{G}}$ that matches a target Gaussian budget $\bar{N}_\mathcal{G}$, we precompute a threshold--budget lookup table, $\tilde{\boldsymbol{\tau}}$ and $\tilde{\mathbf{N}}_\mathcal{G}$, from the predicted multi-level densification score maps $\{\{\hat{\mathbf{D}}_i^l\}_{i=1}^{N_\tec}\}_{l=1}^{L-1}$ using \cref{alg:threshold_budget_lookup}. 


\begin{algorithm}[ht]
\caption{\textbf{Precomputing the threshold--budget lookup table}}
\label{alg:threshold_budget_lookup}
\begin{algorithmic}[1]
% {\hat{\mathbf{D}}_i^l\}_{i=1,l=1}^{N_\text{ctx},L-1}$ 
\Require Multi-level score maps $\{\{\hat{\mathbf{D}}_i^l\}_{i=1}^{N_\text{ctx}}\}^{L-1}_{l=1}$ for all views and levels
\Ensure A sorted threshold list $\tilde{\boldsymbol{\tau}}$ and the corresponding Gaussian-count delta list $\tilde{\mathbf{N}}_\mathcal{G}$

\State Initialize all $\Delta N_i^l$ to $3$ for all values with the same resolution as $\hat{\mathbf{D}}_i^l$.
\For{$l=L-1,L-2,\dots,2$}
    % \State Compute the auto-split mask
    \State $\{\mathbf{A}_i^l\}_{i=1}^{N_\text{ctx}} \leftarrow \left\{ \mathds{1}_{\{\hat{\mathbf{D}}_i^l \ge \mathrm{Up}(\hat{\mathbf{D}}_i^{l-1};2)\}}\right\}_{i=1}^{N_\text{ctx}}$
    % \Comment{automatically densified region when densifying at the level below}
    
    \State $\{\Delta N_i^{l-1} \}_{i=1}^{N_\text{ctx}} \leftarrow \left\{ \Delta N_i^{l-1} + \mathrm{SumPool}_{2\times2}\!\left(\Delta N_i^{l} \odot \mathbf{A}_i^l\right)\right\}_{i=1}^{N_\text{ctx}}$
    \State $\{\Delta N_i^{l} \}_{i=1}^{N_\text{ctx}} \leftarrow \left\{ \Delta N_i^{l} \odot (\mathbf{1}-\mathbf{A}_i^l) \right\}_{i=1}^{N_\text{ctx}}$
\EndFor

\State $\boldsymbol{{\tau}} = \mathrm{Concat}\!\left(\left\{\mathrm{Flatten}(\hat{\mathbf{D}}_i^l)\right\}_{i=1,\dots,N_{\tec};\,l=1,\dots,L-1}\right)$
\State $\boldsymbol{\Delta}\mathbf{N}_\mathcal{G} = \mathrm{Concat}\!\left(\left\{\mathrm{Flatten}(\Delta N_i^l)\right\}_{i=1,\dots,N_{\tec};\,l=1,\dots,L-1}\right)$

\State $\boldsymbol{\pi} = \text{argsort}(\boldsymbol{{\tau}})$
\Comment{descending sorting permutation}

\State $\tilde{\boldsymbol{{\tau}}} = \boldsymbol{{\tau}}[\boldsymbol{\pi}]$
% \Comment{reordering by permutation $\boldsymbol{\pi}$}

\State ${\boldsymbol{\Delta} \mathbf{N}}_\mathcal{G} = \boldsymbol{\Delta}\mathbf{N}_\mathcal{G}[\boldsymbol{\pi}]$
% \Comment{reordering by permutation $\boldsymbol{\pi}$}

\State $\tilde{\mathbf{N}}_\mathcal{G} =  N_{\text{ctx}}H^1W^1 \cdot \mathbf{1} + \text{Cumsum}\left( {\boldsymbol{\Delta} \mathbf{N}}_\mathcal{G} \right) $
\Comment{sorted in ascending order, unlike $\tilde{\boldsymbol{\tau}}$ }

\State \Return $\tilde{\boldsymbol{{\tau}}}, \tilde{\mathbf{N}}_\mathcal{G}$
\end{algorithmic}
\end{algorithm}


Let $\tilde{\mathbf{N}}_\mathcal{G} = \big([\tilde{\mathbf{N}}_\mathcal{G}]_1,\dots,[\tilde{\mathbf{N}}_\mathcal{G}]_K\big)$
denote the Gaussian-count lookup table sorted according to the threshold list
$\tilde{\boldsymbol{\tau}}$. Under the assumption that all densification score values are unique,
the difference between two adjacent entries is upper-bounded by $4^{L-1}-1$:
\[
[\tilde{\mathbf{N}}_\mathcal{G}]_{k+1} - [\tilde{\mathbf{N}}_\mathcal{G}]_{k} \le 4^{L-1}-1.
\]
This follows because lowering the threshold activates exactly one new score value, which changes the allocation at only one spatial position. The largest possible increase occurs when the activated position corresponds to the coarsest valid level and represents all unresolved descendants across finer levels. Since each finest-level position contributes $3$ Gaussians, the maximum increment becomes
\[
3\sum_{j=0}^{L-2}4^j = 4^{L-1}-1.
\]
Therefore, for any target budget $\bar{N}_\mathcal{G}$, there always exists an index $k^\star$ such that
\[
[\tilde{\mathbf{N}}_{\mathcal{G}}]_{k^\star} \le \bar{N}_\mathcal{G}
< [\tilde{\mathbf{N}}_{\mathcal{G}}]_{k^\star} + (4^{L-1}-1),
\]
which implies that the threshold $\tau_{\bar{N}_\mathcal{G}}=\tilde{\boldsymbol{\tau}}_{k^\star}$ satisfies~\cref{eq:budget_condition} of the main paper.
Ultimately, we can efficiently obtain the threshold $\tau_{\bar{N}_\mathcal{G}}$ via \cref{alg:find_budget_threshold}.
% Algorithm 2
\begin{algorithm}[ht]
\caption{\textbf{Finding the budget-matching threshold}}
\label{alg:find_budget_threshold}
\begin{algorithmic}[1]
\Require $\tilde{\boldsymbol{\tau}}$, $\tilde{\mathbf{N}}_{\mathcal{G}}$, $\bar{N}_{\mathcal{G}}$
\Ensure $\tau_{\bar{N}_{\mathcal{G}}}$ satisfying the target budget $\bar{N}_{\mathcal{G}}$

\State $k^\star \leftarrow \max \left\{ k \mid [\tilde{\mathbf{N}}_{\mathcal{G}}]_{k} \le \bar{N}_{\mathcal{G}} \right\}$ 
\Comment{found by binary search in $O(\log |\tilde{\mathbf{N}}_{\mathcal{G}}|)$ time}

\State $\tau_{\bar{N}_{\mathcal{G}}} \leftarrow [\tilde{\boldsymbol{\tau}}]_{k^\star}$
\State \Return $\tau_{\bar{N}_{\mathcal{G}}}$
\end{algorithmic}
\end{algorithm}